\section{Postman, and what introspection costs}
\label{sec:ch02-postman}
\index{Postman|(}

Nitro is a good IDE and it is attached to one service. From
chapter~\ref{ch:enter-the-router} onward Mosaic is several services behind a
router, and the questions worth asking are about the traffic between them: what
the router sent a subgraph, what came back, whether a header survived the trip.
Those are Postman questions, and it is easier to build the habit now, on one
service, than to learn the tool and federation in the same week.

Postman's GraphQL request type reads a schema three ways: by introspecting a
running endpoint, by pulling an API out of a Postman workspace, or from a local
file or URL~\autocite{postman2026client}. Point it at
\code{http://localhost:5100/graphql}, let it introspect, and the query editor
starts autocompleting Mosaic's fields.

The companion repository ships a collection and an environment under
\code{postman/}, with five requests in it. Import both, and the environment
supplies \code{baseUrl} so nothing is hard-coded to a port. Every request
carries assertions, which means the collection doubles as a smoke test: the
product-page request checks that a field from each of four domains came back,
and the whole collection runs from the command line.

\subsection*{Variables, and why to bother early}

The second request in the collection asks for a product page. The identifier is
a variable rather than a literal:

\begin{minted}{graphql}
query ProductPage($id: ID!) {
  productById(id: $id) {
    sku
    title
    price { amount currency }
    availableQuantity
    averageRating
    reviews {
      rating
      author { displayName }
    }
  }
}
\end{minted}

\begin{minted}{json}
{ "id": "a0000000-0000-4000-8000-000000000001" }
\end{minted}

Editing the Variables pane rather than the query is a small discipline with a
large payoff later. A document that stays constant across runs is a document
that can be registered ahead of time, which is what persisted operations and
trusted documents need, and chapter~\ref{ch:the-wire-completely} is largely
about that. Getting into the habit while the graph has one service is free.

\subsection*{A deliberate 400}

\index{Postman!error responses}%
The fourth request in the collection asks for a field that does not exist. I
left it in deliberately, because the response teaches something:

\begin{minted}[fontsize=\footnotesize]{json}
{
  "errors": [
    {
      "message": "The field `nope` does not exist on the type `Product`.",
      "locations": [ { "line": 1, "column": 14 } ],
      "extensions": {
        "type": "Product",
        "field": "nope",
        "responseName": "nope",
        "specifiedBy": "https://spec.graphql.org/September2025/#sec-Field-Selections"
      }
    }
  ]
}
\end{minted}

The status is 400 and there is no \code{data} key at all, not even a null one.
Validation runs against the schema before a single resolver is called, so
nothing was executed and there is nothing partial to return. Notice also what
the error carries: the type, the field, and a link into the GraphQL
specification. There is no HotChocolate error code, because this is not
HotChocolate's rule being broken, it is the language's. Ask for a code and you
find one only on the errors that belong to the server: a syntax error comes back
with \code{extensions.code} set to \code{HC0011}. That division is a useful
thing to have internalised before chapter~\ref{ch:composition}, where reading
error payloads carefully stops being optional.

The two 400s look alike from Postman and they are produced in different places.
An unknown field is a validation rule running inside the request pipeline,
against the schema, on a document that parsed. A syntax error never gets that
far: the transport parses the document before a request exists at all, and what
comes back is the parser refusing to build one.
Chapter~\ref{ch:the-life-of-a-request} takes the pipeline apart and shows where
each of the two happens.

\subsection*{Introspection is off in production, and it will catch you}

\index{introspection!disabled by default}%
This one cost me an hour on this chapter's own code.
\code{AddGraphQLServer()} applies a set of defaults, and one of them disables
introspection whenever the host environment is not Development. The reasoning
is sound: an introspectable schema is a complete map of your API for anybody who
asks, and the safe default is to hand it out only to developers.

The trap is that you can leave Development without meaning to. Mosaic's
verification script starts the service with \code{--no-launch-profile}, which
skips \code{launchSettings.json}, and \code{launchSettings.json} was the only
thing setting the environment. So the service came up in Production, four of the
five Postman requests passed, and the introspection request failed with a 400
and error \code{HC0046}. Nothing about the query was wrong. The fix is to say
what you mean:

\begin{minted}{text}
ASPNETCORE_ENVIRONMENT=Development
\end{minted}

The same applies in a container, and the compose file in the companion
repository sets it with a comment explaining that this is a development
baseline rather than a production posture.

\index{introspection!DisableIntrospection}%
While we are here: if you go looking for how to control this deliberately, the
version~16 security documentation tells you to call
\code{AllowIntrospection(false)} on the builder. That will not
compile~\autocite{chillicream2026introspection}. No such method exists on
\code{IRequestExecutorBuilder} in 16.6.0. \code{AllowIntrospection} lives on
\code{OperationRequestBuilder}, where it overrides the policy for a single
request from inside an interceptor, which is a genuinely useful thing but not
what the page is describing. The builder-level API is \code{DisableIntrospection},
which takes either a boolean or a function of the service provider, and the
function form is how the default is implemented. I report this not to score a
point against a documentation page but because it is exactly the kind of thing
you will blame on your own setup, and you should not have to.

\subsection*{Running the collection without Postman}

The collection is a file, so it runs anywhere:

\begin{minted}{text}
npm ci
npx newman run postman/mosaic.postman_collection.json \
    -e postman/mosaic.local.postman_environment.json
\end{minted}

Five requests, eighteen assertions. Mosaic's verify script runs exactly this
after starting the service, which is what makes the collection a gate rather
than a convenience. Chapter~\ref{ch:the-postman-federation-workbook} builds this
out properly, with JWT pre-request scripts, subscription testing and a
collection that exercises a federated graph through the router. For now the
useful idea is smaller: the requests you use to poke at a service by hand are
worth keeping, and once they have assertions on them they are tests.

Run the collection and Mosaic passes every assertion in it. That is worth
knowing before the next section, which is about how much work it did to get
there.
\index{Postman|)}
